\section{Expected contribution}
The project will deliver a review paper and an assessment matrix that links
requirements to candidate metrics, their prerequisites and their limitations.
The matrix will distinguish measurements of images and datasets from
measurements of a trained model on real data. Recommendations will be
conditional on the vision task, available evidence and resources.
No experimental validation or universally valid quality threshold is promised.

\section{Work plan and collaboration}
Table~\ref{tab:proposal-plan} aligns the work with the confirmed course
deadlines in 2026. Both authors will contribute to screening, verify one another's
extractions and review each other's writing.
\begin{table}[h]
\small
\renewcommand{\arraystretch}{1.2}
\begin{tabularx}{\linewidth}{@{}p{26mm}YY@{}}
\toprule
\textbf{Date (2026)} & \textbf{Activity} & \textbf{Output / milestone} \\
\midrule
22 Sep--4 Oct & Pilot searches, screen, extract and draft initial synthesis. & Documented search and initial evidence matrix. \\
5 Oct & Hand in literature-review draft. & Draft review submission. \\
6--9 Oct & Refine gap, questions, method and feasibility. & Proposal due 9 Oct, 17:00. \\
12 Oct & Participate in peer review. & Review comments and response log. \\
13--18 Oct & Address comments; strengthen evidence. & Revised review and proposal. \\
19 Oct & Receive supervisor feedback. & Agreed revision priorities. \\
20 Oct--5 Nov & Fill evidence gaps; finalise synthesis; proofread. & Checked final documents. \\
6 Nov & Submit both final documents. & Proposal and review due 17:00. \\
\bottomrule
\end{tabularx}
\caption{Course milestones and planned work; timed deadlines use local course time.}
\label{tab:proposal-plan}
\end{table}

The repository will hold separate source files for the paper, proposal,
slides and study records. Changes will be made on focused Git branches and
reviewed by the other author through pull requests. A shared bibliography
and canonical question definitions will keep the documents consistent.
Responsibility for individual sections will be agreed in repository issues;
both authors remain responsible for the integrated argument.

\section{Feasibility, risks and responsible practice}
The review requires literature access, a reference manager, a shared Git
repository and a LaTeX installation. Real datasets, pretrained networks and
training compute are objects of investigation, not prerequisites for carrying
out this literature review. If a later empirical extension is proposed, its
data access, compute and evaluation protocol will need separate planning.

The short period before the first draft makes scope control essential.
We will pilot the search immediately and narrow the application scope with
supervisor agreement if screening is infeasible, documenting the change.
The first draft will distinguish completed synthesis from remaining work.
The main risks are excessive scope, heterogeneous evaluation protocols and
limited evidence for some dimensions. We will pilot the search, document any
scope changes and report evidence gaps rather than force comparisons.
We will distinguish peer-reviewed work from preprints and acknowledge language,
date-range and publication biases. The planned review collects no participant
data and trains no models. We will respect source licences and institutional
access conditions, and keep concise evidence notes with traceable citations.
Before submission we will confirm applicable course and institutional ethics
requirements with the supervisor; no approval or exemption is claimed here.
Human-participant assessment or sensitive-data analysis would require a new
protocol and an institutional ethics determination before starting.
Both authors will verify citations and follow course rules for disclosing
AI-assisted work; no sources or findings will be fabricated.
